\documentclass[12pt]{article}
\usepackage{sbc-template}
\usepackage{graphicx,url}
\usepackage[utf8]{inputenc}
\usepackage[brazilian]{babel}
     
\sloppy

\title{Taxonomia das Arquiteturas de Memória Cache:\\ Estudo Comparativo entre Métodos de
Mapeamento e Políticas de Escrita}

\author{Filipe L. Ribeiro, Davi Erlon L. Morais}


\address{Colegiado de Ciência da Computação -- Universidade Federal do Vale do São Franciso
  (Univasf)\\
  56.304-917 -- Petrolina -- PE -- Brazil
  \email{\{filipe.leiter,davi.erlon\}@discente.univasf.edu.br}
}

\begin{document} 

\maketitle

\begin{abstract}
  This meta-paperilooo describes the style to be used in articles and short papers
  for SBC conferences. For papers in English, you should add just an abstract
  while for the papers in Portuguese, we also ask for an abstract in
  Portuguese (``resumo''). In both cases, abstracts should not have more than
  10 lines and must be in the first page of the paper.
\end{abstract}
     
\begin{resumo} 
  [A taxonomia das arquiteturas de memória cache]
  [Mapeamentos já estudados(direto, totalmente associativo e associativo por conjunto)]
  [Políticas de escrita (duas dimensões ortogonais - acerto E erro)] 
\end{resumo}

\section{Introdução}

A distância entre a velocidade dos processadores e a latência da memória principal desafia a 
arquitetura de computadores há décadas. Em 1965, Moore observou que o número de componentes por 
circuito integrado dobrava a cada ano, projeção que se converteu, ao longo das gerações seguintes, 
em ganho exponencial de desempenho dos processadores (MOORE, 1965). A memória principal, 
mais lenta, não acompanhou esse ritmo, e o desequilíbrio crescente entre as duas curvas foi 
denominado memory wall por Wulf e McKee (1995). A memória cache consolidou-se, então, como 
solução estrutural para esse problema, ao se inserir entre o processador e a memória principal uma 
hierarquia de memórias pequenas e rápidas, em que explora-se a localidade temporal e espacial dos 
programas e evita-se que cada acesso custe toda a latência da DRAM (TANENBAUM; AUSTIN, 2013).

Tema do paragrafo:
O mapeamento, que define como os blocos da memória principal vão se distribuir pelas
linhas da cache, através de diferentes métodos(são muitos, quais falar?)
A politica de substituição que decide qual linha da cache é descartada quando um conjunto se 
completa; A política de escrita, que determina quando e como as modificações alcançam a memória 
principal(Smith, 1982) [Falar sobre problemas que esses entes trazem] \textbf{one}-page texts.

Tema do paragráfo:
Otimizar o projeto de uma memória cache geralmente envolve quatro aspectos:
Maximizar a probabilidade de encontrar na cache o alvo de uma referência à memória (a taxa de acerto);
Minimizar o tempo para acessar informação que de fato está na cache (tempo de acesso);
Minimizar o atraso devido a uma falta; e
Minimizar as sobrecargas de atualização da memória principal, de manutenção da consistência entre 
múltiplas caches etc.

\section{Metodologia de pesquisa}

\section{Taxonomia e Sistematização Teórica}

\section{Análise Crítica e Gargalos Arquiteturais}

\section{Considerações Finais}

\end{document}
